\section{Candidate Solutions}
\label{sec:concepts}

\begin{guidance}[What this section does -- about 200 words plus sketches]
Sketches and short descriptions of several candidate solutions, including
the possible enabling technologies or systems. A single concept is usually a
weak point. Keep the concepts at the same level of detail so that the
comparison in \cref{sec:options} is fair. Methods for generating concepts:
\textcite{cross2008}, \textcite{dym2009}.
\end{guidance}

\begin{figure}[htbp]
  \centering
  \begin{subfigure}[b]{0.32\linewidth}
    \placeholder{Concept A sketch}
    \caption{Concept A}
  \end{subfigure}\hfill
  \begin{subfigure}[b]{0.32\linewidth}
    \placeholder{Concept B sketch}
    \caption{Concept B}
  \end{subfigure}\hfill
  \begin{subfigure}[b]{0.32\linewidth}
    \placeholder{Concept C sketch}
    \caption{Concept C}
  \end{subfigure}
  \caption{Candidate solutions (replace the placeholders with your sketches)}
  \label{fig:concepts}
\end{figure}
